%% =============================================================================
%% Projeto Base - Relatório de Estágio com abnTeX2
%% Customizado para a UNIJUÍ
%% Template desenvolvido e mantido por: Prof. Eldair F. Dornelles
%% A versão mais atualizada deste template pode ser obtida no GitHub:
%% https://github.com/dornelles/relatorio-de-estagio-template-latex
%% Qualquer sugestão de melhoria será bem vinda, 
%% Contato pelo e-mail: eldair.dornelles@unijui.edu.br
%% =============================================================================

\documentclass[
    12pt,               % tamanho da fonte
    oneside,            % para impressão frente e verso (use oneside para apenas frente)
    a4paper,            % tamanho do papel
    chapter=TITLE,      % títulos de capítulos em maiúsculas
    section=TITLE,      % títulos de seções em maiúsculas
    english,            % idioma adicional para hifenização
    brazil              % o último idioma é o principal do documento
]{abntex2}

% ==============================================================================
% PACOTES ESSENCIAIS
% ==============================================================================

% --- Codificação e fontes ---
\usepackage[T1]{fontenc}        % seleção de códigos de fonte
\usepackage[utf8]{inputenc}     % codificação do documento (conversão automática dos acentos)
% \usepackage{lmodern}          % comentado: a fonte agora está definida no arquivo unijui.sty

% --- Pacotes de citações ---
\usepackage[brazilian,hyperpageref]{backref}  % páginas com as citações na bibl
\usepackage[alf,bibjustif]{abntex2cite}        % citações padrão ABNT (bibjustif = referências justificadas)

% --- Pacotes matemáticos ---
\usepackage{amsmath}
\usepackage{amssymb}

% --- Pacotes gráficos ---
\usepackage{graphicx}           % inclusão de gráficos
% NOTA: NÃO usar \usepackage{float} com abnTeX2/memoir - conflita com \newfloat
\usepackage{subfig}             % subfiguras
\usepackage{pdfpages}           % inclusão de páginas PDF

% --- Pacotes de tabelas ---
\usepackage{booktabs}           % tabelas profissionais
\usepackage{multirow}           % células que ocupam várias linhas
\usepackage{longtable}          % tabelas longas (múltiplas páginas)
\usepackage{array}              % personalização de colunas

% --- Outros pacotes úteis ---
\usepackage{microtype}          % melhorias de justificação
\usepackage{indentfirst}        % indenta o primeiro parágrafo de cada seção
\usepackage{color}              % controle de cores
\usepackage{xcolor}             % cores avançadas
\usepackage{listings}           % listagens de código-fonte
\usepackage{enumitem}           % listas personalizadas
\usepackage{lipsum}             % geração de texto fictício (remover na versão final)
\usepackage{totcount}           % conta o total de floats (figuras, tabelas, etc.) no documento

% --- Customizações UNIJUÍ ---
\input{unijui.sty}

% ==============================================================================
% CONFIGURAÇÕES DO PACOTE BACKREF
% ==============================================================================
% Configurações do pacote backref
% Usado sem a opção hyperpageref de backref
\renewcommand{\backrefpagesname}{Citado na(s) página(s):~}
% Texto padrão antes do número das páginas
\renewcommand{\backref}{}
% Define os textos da citação
\renewcommand*{\backrefalt}[4]{
    \ifcase #1 %
        Nenhuma citação no texto.%
    \or
        Citado na página #2.%
    \else
        Citado #1 vezes nas páginas #2.%
    \fi}%

% ==============================================================================
% CONFIGURAÇÕES DE LISTAGENS DE CÓDIGO
% ==============================================================================
\lstset{
    basicstyle=\ttfamily\small,
    numbers=left,
    numberstyle=\tiny\color{gray},
    stepnumber=1,
    numbersep=8pt,
    backgroundcolor=\color{white},
    showspaces=false,
    showstringspaces=false,
    showtabs=false,
    frame=single,
    rulecolor=\color{black},
    tabsize=2,
    captionpos=b,
    breaklines=true,
    breakatwhitespace=false,
    keywordstyle=\color{blue}\bfseries,
    commentstyle=\color{green!60!black},
    stringstyle=\color{red},
    escapeinside={\%*}{*)},
    literate=
        {á}{{\'a}}1 {é}{{\'e}}1 {í}{{\'i}}1 {ó}{{\'o}}1 {ú}{{\'u}}1
        {Á}{{\'A}}1 {É}{{\'E}}1 {Í}{{\'I}}1 {Ó}{{\'O}}1 {Ú}{{\'U}}1
        {ã}{{\~a}}1 {õ}{{\~o}}1 {Ã}{{\~A}}1 {Õ}{{\~O}}1
        {ç}{{\c{c}}}1 {Ç}{{\c{C}}}1
        {ê}{{\^e}}1 {â}{{\^a}}1 {ô}{{\^o}}1
}

% ==============================================================================
% DADOS DO TRABALHO
% ==============================================================================

% --- Informações do documento (preencher conforme o trabalho) ---
\titulo{Título do Relatório}
\autor{Nome do Estudante}
\local{Ijuí/RS}
\data{ANO DA ENTREGA}
\orientador{Prof. xxxx}
\instituicao{%
    Universidade Regional do Noroeste do\\[1mm]
    Estado do Rio Grande do Sul
}

% --- Informações específicas da UNIJUÍ (definidas no unijui.sty) ---
\campus{Ijuí}
\curso{Ciência da Computação}
\professordabanca{Prof(a). xxxx}
\datadabanca{XX}{XXX}{XXXX}

% --- Preâmbulo (texto da folha de rosto e aprovação) ---
% Não adicione ponto final do texto. Na folha de rosto o ponto irá aparecer automaticamente, 
% já na folha de aprovação ao invés de ponto irá aparecer uma vírgula, seguida do texto 
% "submetido à aprovação da banca examinadora composta pelos seguintes membros:"
\preambulo{Relatório apresentado ao curso de xxxxx da Universidade Regional do Noroeste do Estado do Rio Grande do Sul - UNIJUÍ, como requisito para a obtenção do grau de xxxx}

% --- Resumo e Abstract ---
\resumoemportugues{%
    Escreva aqui o resumo do trabalho em português. O resumo deve conter
    entre 150 e 500 palavras e apresentar de forma concisa os objetivos,
    a metodologia, os resultados e as conclusões do trabalho.
}
% Instrução: Use letras minúsculas para as palavras-chave, separe-as com ponto e vírgula (;) e encerre a última com um ponto (.)
\palavraschaveemportugues{palavra-chave-1; palavra-chave-2; palavra-chave-3.}


% ==============================================================================
% CONFIGURAÇÕES DE APARÊNCIA DO PDF FINAL
% ==============================================================================
\makeatletter
\hypersetup{
    pdftitle={\@title},
    pdfauthor={\@author},
    pdfsubject={\imprimirpreambulo},
    pdfcreator={LaTeX com abnTeX2},
    pdfkeywords={abnt}{latex}{abntex}{abntex2}{relatório de estágio},
    colorlinks=true,        % false: links em caixas; true: links coloridos
    linkcolor=black,        % cor dos links internos
    citecolor=black,        % cor dos links para bibliografias
    filecolor=black,        % cor dos links de arquivos
    urlcolor=black,         % cor dos links de URLs
    bookmarksdepth=4
}
\makeatother

% ==============================================================================
% ESPAÇAMENTOS ENTRE LINHAS E PARÁGRAFOS
% ==============================================================================
% O tamanho do parágrafo é dado por:
\setlength{\parindent}{1.3cm}
% Controle do espaçamento entre um parágrafo e outro:
\setlength{\parskip}{0.2cm}  % tente também \onelineskip

% ==============================================================================
% INÍCIO DO DOCUMENTO
% ==============================================================================
\begin{document}

% Seleciona o idioma do documento
\selectlanguage{brazil}

% Retira espaço extra obsoleto entre as frases
\frenchspacing

% --- Imprimir página de créditos e marca d'água ---
% \imprimirpaginadecreditos
% \imprimirmarcadagua

% ==========================================================================
% ELEMENTOS PRÉ-TEXTUAIS
% ==========================================================================
\pretextual

% --- Capa ---
\imprimircapa

% --- Folha de rosto ---
% (o * indica que haverá a ficha bibliográfica)
\imprimirfolhaderosto

% --- Ficha catalográfica (inserir PDF gerado pela biblioteca) ---
% \begin{fichacatalografica}
%     \includepdf{fig_ficha_catalografica.pdf}
% \end{fichacatalografica}

% --- Folha da banca ---
\imprimirfolhadabanca

% --- Agradecimentos ---
\begin{agradecimentos}
Agradeço a todos que contribuíram para a realização deste trabalho.

\end{agradecimentos}

% % --- Epígrafe ---
% \begin{epigrafe}
%     \vspace*{\fill}
%     \begin{flushright}
%         \textit{``A imaginação é mais importante que o conhecimento.''\\
%         (Albert Einstein)}
%     \end{flushright}
% \end{epigrafe}

% --- Resumo em português ---
\begin{resumo}
    \imprimirresumo
    \par
    \noindent \textbf{Palavras-chave}: \imprimirpalavraschave
\end{resumo}

% --- Lista de ilustrações (só aparece se houver figuras) ---
% Na primeira compilação os contadores podem ser 0; compile 2x para resultado correto.
\regtotcounter{figure}
\ifnum\totvalue{figure}>0
    \pdfbookmark[0]{\listfigurename}{lof}
    \listoffigures*
    \cleardoublepage
\fi

% --- Lista de quadros (só aparece se houver quadros) ---
\regtotcounter{quadro}
\ifnum\totvalue{quadro}>0
    \pdfbookmark[0]{\listofquadrosname}{loq}
    \listofquadros*
    \cleardoublepage
\fi

% --- Lista de tabelas (só aparece se houver tabelas) ---
\regtotcounter{table}
\ifnum\totvalue{table}>0
    \pdfbookmark[0]{\listtablename}{lot}
    \listoftables*
    \cleardoublepage
\fi

% --- Lista de abreviaturas e siglas ---
\begin{siglas}
    \item[ABNT] Associação Brasileira de Normas Técnicas
    \item[TCC] Trabalho de Conclusão de Curso
    \item[UNIJUÍ] Universidade Regional do Noroeste do Estado do Rio Grande do Sul
\end{siglas}

% --- Lista de símbolos ---
\begin{simbolos}
    \item[$ \Gamma $] Letra grega Gama
    \item[$ \Lambda $] Lambda
\end{simbolos}

% --- Sumário ---
\pdfbookmark[0]{\contentsname}{toc}
\tableofcontents*
\cleardoublepage

% ==========================================================================
% ELEMENTOS TEXTUAIS
% ==========================================================================
\textual

% --- Inclusão dos capítulos ---
\input{capitulos/01-introducao}
\input{capitulos/02-apresentacao-da-empresa}
\input{capitulos/03-metodologia}
\input{capitulos/04-revisao-bibliografica}
\input{capitulos/05-desenvolvimento}
\input{capitulos/06-consideracoes-finais}

% ==========================================================================
% ELEMENTOS PÓS-TEXTUAIS
% ==========================================================================
\postextual

% --- Referências bibliográficas ---
\bibliography{referencias}

% --- Apêndices ---
\begin{apendicesenv}
    % %Imprime uma página indicando o início dos apêndices
    % \partapendices

    \input{apendices/apendice-a}
\end{apendicesenv}

% --- Anexos ---
\begin{anexosenv}
    % %Imprime uma página indicando o início dos anexos
    % \partanexos

    \input{anexos/anexo-a}
\end{anexosenv}

% --- Índice remissivo (opcional) ---
% \phantompart
% \printindex

\end{document}
